% Compact verifier polytope proof resource.
% This file is intentionally standalone and is also imported conceptually by the IEEE-style note.
\documentclass[11pt]{article}
\usepackage{amsmath,amssymb,amsthm,mathtools,bm}
\usepackage[margin=1in]{geometry}
\usepackage{booktabs}
\newtheorem{proposition}{Proposition}
\newtheorem{lemma}{Lemma}
\newtheorem{remark}{Remark}
\newcommand{\R}{\mathbb{R}}
\newcommand{\norm}[1]{\left\lVert #1\right\rVert_2}
\newcommand{\trans}{^{\mathsf T}}
\title{Compact Verifier Polytope: SOCP Proof Notes}
\author{Safe Flow Expansion / Pillar 3}
\date{}
\begin{document}
\maketitle

\section{Local Level-Set Certificate}
Let $c=q_0\in\R^2$ and let $q_0,q_1,\ldots,q_H$ be the local MPPI rollout.  Define
\[
    \alpha_t=(1-\gamma)^t,\qquad \beta_t=1-\alpha_t.
\]
A face is represented as
\[
    a_i\trans(x-c)\le m_i.
\]
Its normalized face barrier is
\[
    h_i(x)=\frac{m_i-a_i\trans(x-c)}{m_i}.
\]
The condition $h_i(q_t)\ge\alpha_t$ is equivalent to
\[
    a_i\trans(q_t-c)\le \beta_t m_i.
\]
\textbf{Thus the recursive level-set certificate is linear in $(a_i,m_i)$.}

\section{Circle Separation as a Cone Constraint}
For disk $O_i=\{x:\norm{x-o_i}\le r_i\}$, the face separates the disk if
\[
    a_i\trans(x-c)\ge m_i\qquad\forall x\in O_i.
\]
The disk support identity is
\[
    \min_{x\in O_i}a_i\trans(x-c)=a_i\trans(o_i-c)-r_i\norm{a_i}.
\]
Therefore disk separation is
\[
    r_i\norm{a_i}\le a_i\trans(o_i-c)-m_i.
\]
\textbf{This is exactly a second-order cone constraint.}

\section{Margin-Bounded SOCP}
The compact verifier solves
\begin{align}
    \max_{\{a_i,m_i\}}\quad &\sum_i w_i m_i \\
    \mathrm{s.t.}\quad
    &a_i\trans(q_t-c)\le\beta_t m_i, &&\forall i,t,\\
    &r_i\norm{a_i}\le a_i\trans(o_i-c)-m_i, &&\forall i,\\
    &\norm{a_i}\le1, &&\forall i,\\
    &m_{\min}\le m_i\le m_{\max}, &&\forall i.
\end{align}

\begin{proposition}[SOCP structure]
The margin-bounded compact verifier is an SOCP.
\end{proposition}
\begin{proof}
The trajectory constraints are affine.  The disk-separation inequalities are Lorentz cone inequalities of the form $\norm{Az+b}\le c\trans z+d$.  The bound $\norm{a_i}\le1$ is also SOC.  The margin bounds are affine, and the objective is linear.  Hence the problem is an SOCP.
\end{proof}

\begin{proposition}[Unit-normal relaxation without an upper margin cap]
If $w_i>0$, $m_i>0$, and the upper bound $m_i\le m_{\max}$ is inactive or absent, then an optimum of an independent face saturates $\norm{a_i}=1$.
\end{proposition}
\begin{proof}
For $0<\norm{a_i}<1$, choose $s>1$ such that $s\norm{a_i}\le1$ and set $\tilde a_i=s a_i$, $\tilde m_i=s m_i$.  The trajectory and disk constraints are homogeneous in $(a_i,m_i)$, so feasibility is preserved. Since $w_i>0$, $w_i\tilde m_i>w_i m_i$.  Therefore a non-saturated solution cannot be optimal while the upper margin cap is inactive.
\end{proof}

\begin{remark}[Effect of $m_{\max}$]
\textbf{The upper bound $m_i\le m_{\max}$ is a geometry-changing constraint, not a weight rescaling.}  If the unconstrained max-margin face is too wide or produces a pointy wedge, choosing $m_{\max}$ near the desired tube half-width clips the face and can produce a compact tube-like verifier polytope.
\end{remark}

\section{Artificial Boundary Anchors}
For sensing radius $R$ and $K$ directions
\[
 n_\ell=(\cos(2\pi\ell/K),\sin(2\pi\ell/K))\trans,
\]
let
\[
    M_K=R\cos(\pi/K),\qquad
    o_\ell^{\rm art}=c+(M_K+\rho_{\rm art})n_\ell,
    \qquad r_\ell^{\rm art}=\rho_{\rm art}.
\]
Then the radial tangent margin is $M_K$.  \textbf{Artificial anchors encode the sensing envelope as ordinary circular obstacles in the same SOCP template.}

\section{Positive Weight Invariance for Independent Faces}
If the feasible set decomposes as $\mathcal C=\prod_i\mathcal C_i$, then
\[
\max_{(a_i,m_i)\in\mathcal C_i}\sum_i w_i m_i
=\sum_i\max_{(a_i,m_i)\in\mathcal C_i}w_i m_i.
\]
For $w_i>0$, $\arg\max w_i m_i=\arg\max m_i$.  \textbf{Thus positive weights alone do not reshape independent one-face-per-anchor polytopes.}

\end{document}
